\documentclass[10pt,twocolumn]{article}
\usepackage[T1]{fontenc}
\usepackage[utf8]{inputenc}
\usepackage{lmodern}
\usepackage[margin=1.8cm,top=2.4cm,bottom=2cm,columnsep=0.6cm]{geometry}
\usepackage{fancyhdr}
\usepackage{titlesec}
\usepackage{booktabs}
\usepackage{amsmath,amssymb,amsfonts}
\usepackage{microtype}
\usepackage{xcolor}
\usepackage{mdframed}
\usepackage{parskip}
\usepackage{enumitem}
\usepackage{hyperref}

\definecolor{rulered}{RGB}{180,30,30}
\definecolor{boxgray}{RGB}{245,245,245}
\definecolor{keywordgray}{RGB}{100,100,100}

\pagestyle{fancy}
\fancyhf{}
\fancyhead[L]{\textsc{\small Claude Mag}}
\fancyhead[C]{\small\textit{Machine Learning Research Digest}}
\fancyhead[R]{\small 2026-10-21}
\fancyfoot[C]{\small\thepage}
\renewcommand{\headrulewidth}{0.8pt}

\titleformat{\section}{\normalsize\bfseries\color{rulered}}{}{0pt}{}%
  [\vspace{-4pt}\color{rulered}\rule{\columnwidth}{0.4pt}]
\titlespacing{\section}{0pt}{8pt}{4pt}

\hypersetup{colorlinks=true,urlcolor=rulered,linkcolor=black}

\begin{document}
\setlength{\parindent}{0pt}
\setlength{\parskip}{4pt}

{\small\color{keywordgray}\textbf{Keywords:}
  visual token compression \textbullet{}
  native MLLM \textbullet{}
  Pareto optimization \textbullet{}
  inference efficiency}\par\vspace{4pt}

{\LARGE\bfseries\raggedright
  PACE: Pareto-Adaptive Compression for Efficient Native MLLMs}\par\vspace{4pt}

{\large\itshape\raggedright
  Orthogonal Spatial Saliency Decomposition Achieves Near-Lossless
  Performance at 25--50\% Token Retention with 2$\times$ FLOPs Reduction}\par\vspace{6pt}

{\small\textbf{By} Lei Tian, Chunzheng Zhu, Kai Wu, Yifan Zhang \&
  Haihua Yang}\par\vspace{2pt}

{\small\color{keywordgray}arXiv:2610.24005 \textbullet{}
  NeurIPS 2026}
\par\vspace{2pt}\noindent{\color{rulered}\rule{\columnwidth}{0.6pt}}\vspace{6pt}

Native multimodal large language models process image patches directly,
without a separate vision encoder, but this architectural simplicity comes
at a cost: no encoder-derived saliency signal is available for token
compression. PACE addresses this gap with an orthogonal spatial attention
decomposition and a Pareto-optimal budget allocation head. At 50\% visual
token retention, PACE delivers near-lossless performance on standard benchmarks
and more than $2\times$ FLOPs reduction; at 25\% retention it improves on
chart-understanding tasks while remaining competitive elsewhere. Accepted at
NeurIPS 2026.

\begin{mdframed}[backgroundcolor=boxgray,linecolor=rulered,linewidth=1pt,
    innerleftmargin=6pt,innerrightmargin=6pt,
    innertopmargin=6pt,innerbottommargin=6pt]
\textbf{\small AT A GLANCE}\par\vspace{4pt}
\begin{description}[leftmargin=0pt,labelsep=4pt,itemsep=2pt,parsep=0pt,
    font=\small\bfseries]
  \item[Problem:] ViT-free (native) MLLMs have no encoder attention signal
    to guide token compression; existing methods are inapplicable.
  \item[Why it matters:] Native MLLMs are increasingly popular; without
    compression, visual token sequences create prohibitive attention cost.
  \item[Method:] Decompose attention into orthogonal height and width
    saliency components gated by a positional signal; allocate token
    budgets via Pareto optimization over performance, efficiency, and
    spatial coverage.
  \item[Result:] Near-lossless accuracy at 50\% retention; $>$2$\times$
    FLOPs reduction; improvement on chart benchmarks at 25\% retention.
  \item[Evidence:] Evaluation across multiple native MLLM architectures
    and benchmarks; ablations confirm each component's contribution.
\end{description}
\end{mdframed}

\section{The Big Picture}

The field has converged on two approaches to multimodal LLMs. The dominant
approach uses a separate vision encoder (ViT) to process images into a
compact visual representation before passing it to the language model.
The emerging alternative---the \textit{native} or \textit{ViT-free} MLLM---
processes raw image patches directly in the transformer, unifying modalities
from the start.

Native architectures reduce engineering complexity and enable end-to-end
optimization, but they reintroduce a well-known problem at scale: the number
of visual tokens grows quadratically with image resolution. For a 448$\times$448
image at a patch size of 14, that is 1024 tokens---before any text. Attention
cost is $O(N^2)$; this is already prohibitive at high resolution or in
multi-image settings.

Compressing visual tokens is the natural remedy, but every existing compression
method was designed for ViT-based pipelines, where the encoder's [CLS]
token attention or patch-similarity matrix provides a natural saliency signal.
Native models have neither. PACE is the first method designed specifically
for this setting.

\section{Why Existing Approaches Fall Short}

\textbf{ViT-based compression (ToMe, FastViT, DynamicViT):} All rely on
the encoder's attention map or token similarity computed within the ViT.
These signals do not exist in native architectures; applying these methods
to native models produces near-random compression.

\textbf{Uniform subsampling:} Discards spatial structure uniformly, missing
locally salient regions that are scattered across the image.

\textbf{Greedy attention thresholding:} Retains the tokens with the highest
self-attention scores in the LLM; this greedy approach ignores spatial
coverage and can cluster selections around a single prominent region, missing
the rest of the image.

\section{Core Mechanism}

PACE introduces two components applied sequentially at each layer of the
native MLLM.

\textbf{Orthogonal spatial saliency decomposition:} The attention matrix
$\mathbf{A} \in \mathbb{R}^{N \times N}$ for the visual tokens is factored
into marginals along the height and width dimensions of the image grid:
\[
\mathbf{A} \approx \mathbf{A}_H \otimes \mathbf{A}_W
\]
where $\mathbf{A}_H \in \mathbb{R}^{h \times h}$ and
$\mathbf{A}_W \in \mathbb{R}^{w \times w}$ capture row-wise and column-wise
saliency, with $N = h \times w$. A positional attention signal $p_{ij}$ (a
learned embedding of the token's position relative to attended regions)
gates the combined score.

\textbf{Pareto budget allocation:} Rather than fixing a single compression
ratio for all inputs, PACE frames token budget selection as a Pareto
optimization over three objectives: performance preservation, FLOPs reduction,
and spatial coverage. A lightweight allocation network selects the operating
point at inference time.

\section{Technical Formulation}

The saliency score for token $(i,j)$ is:
\[
s_{ij} = \alpha\,[\mathbf{A}_H]_{ii} + \beta\,[\mathbf{A}_W]_{jj} + \gamma\, p_{ij}
\]
where $(\alpha, \beta, \gamma)$ are learned scalars. The top-$B$ tokens by
score are retained, where $B$ is determined by the allocation network.

The Pareto budget allocation problem is:
\[
\min_{\mathbf{m}}\;\left(\text{PerfDrop}(\mathbf{m}),\;
  -\text{FLOPs\_red}(\mathbf{m}),\;
  -\text{SpatCov}(\mathbf{m})\right)
\]
subject to $\|\mathbf{m}\|_0 = B$, where $\mathbf{m}$ is the binary
selection mask. The Pareto frontier is computed on a held-out calibration
set and parameterised by a small MLP that maps input complexity features
(image entropy, text query length) to a budget $B$ and weighting
$(\alpha, \beta, \gamma)$.

At training time, the allocation MLP is trained to minimise the weighted
Chebyshev distance to the Pareto frontier:
\[
\mathcal{L}_{\text{Pareto}}(\phi) = \mathbb{E}\!\left[
  \max\!\left(\lambda_1\, \Delta\text{Perf},\;
    \lambda_2\, \Delta\text{FLOPs},\;
    \lambda_3\, \Delta\text{Cov}\right)\right]
\]

\section{Learning or Inference Procedure}

\begin{enumerate}[itemsep=2pt,parsep=0pt]
  \item Construct the Pareto frontier by running the model at several
    compression ratios on a calibration set; record performance, FLOPs,
    and spatial coverage.
  \item Train the allocation MLP on the frontier with $\mathcal{L}_{\text{Pareto}}$.
  \item At inference: the allocation MLP selects budget $B$ and weightings;
    saliency scores are computed; the top-$B$ tokens are retained; the
    remaining tokens are pruned before subsequent transformer layers.
  \item PACE is a plug-in module; the base MLLM is not modified or
    fine-tuned.
\end{enumerate}

\section{What the Guarantee Says}

The Pareto formulation provides a formal guarantee: the selected budget
$B$ lies on the Pareto frontier (no other $B' \neq B$ strictly improves
all three objectives simultaneously). The spatial coverage constraint ensures
that selection never degenerates into clustering around a single image region,
bounding the worst-case missed-region probability.

\section{Experimental Findings}

\begin{itemize}[itemsep=2pt,parsep=0pt]
  \item \textbf{At 50\% token retention:} Near-lossless accuracy on VQA,
    spatial reasoning, and multimodal reasoning benchmarks; FLOPs reduced
    by 2.1--2.4$\times$.
  \item \textbf{At 25\% token retention:} Performance improves on chart
    understanding (spatial coverage prevents missing chart regions that
    a greedy method would skip); slight degradation on spatial reasoning.
  \item \textbf{Cross-architecture:} Gains are consistent across multiple
    native MLLM architectures, confirming that PACE is not tuned to a
    single model's attention pattern.
  \item \textbf{Positional signal ablation:} Removing the positional gate
    $p_{ij}$ hurts spatial-reasoning tasks the most, where token position
    encodes geometric relationships critical for the answer.
\end{itemize}

\section{Ablations and Interpretation}

\textbf{Without Pareto allocation (fixed global ratio):} Uniform 50\%
budget performs well on average but degrades on information-dense inputs
(too aggressive) and wastes compute on information-sparse inputs (not
aggressive enough).

\textbf{Without orthogonal decomposition (scalar per-token attention):}
The saliency map is blurry; clusters of redundant background tokens score
similarly to salient foreground tokens.

\textbf{Height-only or width-only saliency:} Misses salient regions
localised in the orthogonal direction (e.g., a horizontal strip of text
is missed by height-only saliency when the strip occupies few rows).

\textbf{Why native MLLMs tolerate compression:} The authors hypothesise that
without ViT pre-processing, the transformer distributes attention more
uniformly in early layers and learns implicit token-aggregation in later
layers; explicit compression makes this implicit process explicit.

\section*{Reference}

Lei Tian, Chunzheng Zhu, Kai Wu, Yifan Zhang, Haihua Yang.
\textbf{PACE: Pareto-Adaptive Compression for Efficient Native MLLMs}.
arXiv:2610.24005, NeurIPS 2026.
\url{https://arxiv.org/abs/2610.24005}

\end{document}
